\editionPage[\pdfbookmark[0]{Indications des défis}{bm-indices-debut-parent}]{indices-debut}{Indications progressives}{Défis 1 à 3}{Ne lire que l’aide nécessaire.}
\begin{blocEdition}{indices-debut}
\editionEtape{D01 · Numéroter tout un livre}\textbf{Indice 1.} Compter séparément les pages à un chiffre, à deux chiffres, à trois chiffres.\par\medskip\textbf{Indice 2.} De 10 à 99, il y a 90 pages ; de 100 à 150, il y en a 51.\par
\editionEtape{D02 · Le carré magique}\textbf{Indice 1.} La somme de trois cases alignées vaut toujours 15 ; que vaut la case opposée à 2 dans la diagonale ?\par\medskip\textbf{Indice 2.} Les nombres pairs 2, 4, 6, 8 occupent les coins.\par
\editionEtape{D03 · Plier, plier, plier}\textbf{Indice 1.} Chaque pli double le nombre d’épaisseurs.\par\medskip\textbf{Indice 2.} Écrire la suite des épaisseurs : 2, 4, 8, …\par

\end{blocEdition}
\editionPage{indices-4}{Indications progressives}{Défis 4 à 6}{Ne lire que l’aide nécessaire.}
\begin{blocEdition}{indices-4}
\editionEtape{D04 · Cent œufs, douze boîtes}\textbf{Indice 1.} Si les 12 boîtes étaient toutes des boîtes de 6, combien d’œufs seraient rangés ?\par\medskip\textbf{Indice 2.} Remplacer une boîte de 6 par une boîte de 10 ajoute 4 œufs.\par
\editionEtape{D05 · La tour triangulaire}\textbf{Indice 1.} Écrire la suite : 1, 3, 6, 10, … en ajoutant à chaque fois l’étage suivant.\par\medskip\textbf{Indice 2.} À l’étape 8, on additionne les nombres de 1 à 8.\par
\editionEtape{D06 · Le drapeau étoilé}\textbf{Indice 1.} Partager chaque bande en deux pour obtenir des parts toutes identiques.\par\medskip\textbf{Indice 2.} Compter les parts du nouveau partage, puis les parts étoilées.\par

\end{blocEdition}
\editionPage{indices-7}{Indications progressives}{Défis 7 à 9}{Ne lire que l’aide nécessaire.}
\begin{blocEdition}{indices-7}
\editionEtape{D07 · La course des fourmis}\textbf{Indice 1.} Comparer d’abord chaque position à 1.\par\medskip\textbf{Indice 2.} Pour départager A et B, les parts sont déjà de même taille.\par
\editionEtape{D08 · Le nombre au code secret}\textbf{Indice 1.} Appeler u le chiffre des unités : que vaut le chiffre des dixièmes ?\par\medskip\textbf{Indice 2.} La somme des deux chiffres fait trois fois u.\par
\editionEtape{D09 · Entre deux graduations}\textbf{Indice 1.} Écrire les bornes en centièmes : 7,20 et 7,30.\par\medskip\textbf{Indice 2.} Lister les centièmes possibles sans prendre les bornes.\par

\end{blocEdition}
\editionPage{indices-10}{Indications progressives}{Défis 10 à 12}{Ne lire que l’aide nécessaire.}
\begin{blocEdition}{indices-10}
\editionEtape{D10 · Cinq euros, huit pièces}\textbf{Indice 1.} Essayer d’abord sans pièce de 2 euros.\par\medskip\textbf{Indice 2.} Les pièces de 50 centimes vont par deux pour faire des euros entiers.\par
\editionEtape{D11 · La corde à border}\textbf{Indice 1.} Le périmètre du carré est le quadruple du côté.\par\medskip\textbf{Indice 2.} Dans un rectangle, largeur et longueur se répètent deux fois chacune.\par
\editionEtape{D12 · Quatre litres exactement}\textbf{Indice 1.} Remplir le bidon de 3 L et le verser dans le grand.\par\medskip\textbf{Indice 2.} Quand le grand bidon est plein, que reste-t-il dans le petit ?\par

\end{blocEdition}
\editionPage{indices-13}{Indications progressives}{Défis 13 à 15}{Ne lire que l’aide nécessaire.}
\begin{blocEdition}{indices-13}
\editionEtape{D13 · Les pièces du carré}\textbf{Indice 1.} Commencer par l’aire du grand triangle.\par\medskip\textbf{Indice 2.} Le reste se partage en deux aires égales.\par
\editionEtape{D14 · L’horloge pressée}\textbf{Indice 1.} Compter les heures réellement écoulées depuis 8 h.\par\medskip\textbf{Indice 2.} Chaque heure écoulée ajoute 3 minutes d’avance : à 20 h réelles, combien affiche-t-elle ?\par
\editionEtape{D15 · Le trésor du jardin}\textbf{Indice 1.} Construire soigneusement, puis mesurer TA et TB.\par\medskip\textbf{Indice 2.} Le pliage le long de la perpendiculaire superpose A et B.\par

\end{blocEdition}
\editionPage{indices-16}{Indications progressives}{Défis 16 à 18}{Ne lire que l’aide nécessaire.}
\begin{blocEdition}{indices-16}
\editionEtape{D16 · Le pliage du carré}\textbf{Indice 1.} Quels côtés du carré restent visibles après le pliage ?\par\medskip\textbf{Indice 2.} Comparer la diagonale à un côté du carré : laquelle est la plus longue ?\par
\editionEtape{D17 · Les mots du miroir}\textbf{Indice 1.} Le miroir retourne le mot : la dernière lettre devient la première.\par\medskip\textbf{Indice 2.} Il faut des lettres symétriques et un mot qui se lit pareil dans les deux sens.\par
\editionEtape{D18 · Le grand cube peint}\textbf{Indice 1.} Les cubes à 3 faces peintes sont aux coins du grand cube.\par\medskip\textbf{Indice 2.} Un seul petit cube ne touche aucune face extérieure.\par

\end{blocEdition}
\editionPage{indices-19}{Indications progressives}{Défis 19 à 21}{Ne lire que l’aide nécessaire.}
\begin{blocEdition}{indices-19}
\editionEtape{D19 · La photo agrandie}\textbf{Indice 1.} Chercher par combien le petit côté a été multiplié.\par\medskip\textbf{Indice 2.} Le grand côté est multiplié par le même nombre.\par
\editionEtape{D20 · Le tableau effacé}\textbf{Indice 1.} Compléter d’abord la ligne CM2 à l’aide du total général.\par\medskip\textbf{Indice 2.} Dans une ligne ou une colonne, la case manquante se déduit du total.\par
\editionEtape{D21 · Le sac mystère}\textbf{Indice 1.} « Probable » signifie plus de la moitié des boules.\par\medskip\textbf{Indice 2.} « Possible » signifie au moins une boule bleue.\par

\end{blocEdition}
